\documentclass[11pt]{letter}
\usepackage{mathptmx}
\textheight  8.5in
\textwidth   6.in
\voffset     -.5in
\hoffset     -.5in

\newcommand{\ie}  {\textit{i.e., }}
\newcommand{\eg}  {\textit{e.g., }}
\newcommand{\etal}{\textit{et al., }}
\newcommand{\ca}  [1]{\mathcal{#1}}
\newcommand{\bs}  [1]{\mathbf{#1}}

\signature{K. Shafer Smith \\ Geoffrey K. Vallis}
\address{K. Shafer Smith\\  Geoffrey K. Vallis \\
  GFDL/Princeton University\\
  PO Box 308\\
  Princeton, NJ 08542\\
  Tel: (609) 452-6537\\
  Email: kss@gfdl.noaa.gov
}


\begin{document}

\begin{letter}{Dr. William R. Young \\
 Editor, JFM \\
 Physical Oceanography Research Division \\
 Scripps Institution of Oceanography \\
 University of California San Diego \\
 La Jolla, California 92093-0213
 }

\opening{Dear Bill,}

Enclosed please find two copies of a slightly revised version of the
manuscript ``Turbulent Diffusion in the Geostrophic Inverse Cascade.''
We respond to the comments of reviewer B below.

\begin{quote}
p. 4, end of top paragraph:  ``... - then the effects of friction become
similar to [(a) - this is a typo, I suppose?] those of a barrier...'' - I
am not sure what this statement means; it should either be better
explained or dropped.
\end{quote}

This statement was dropped.

\begin{quote}
p. 13, the discussion of the ``halting scale.''  I still don't
understand what this ``halting scale'' means.  If the authors imply
termination of the inverse cascade, then this is a wrong statement;
the rate of the inverse cascade is $\epsilon$ and termination of the
inverse cascade would mean $\epsilon=0$ which the authors themselves
find unphysical.  Equation (5.21) as a definition of the halting scale
does not provide more clarity because epsilon remains $>0$ for wave
numbers smaller than $k_r$ allowing the inverse cascade to penetrate
beyond this ``halting scale.''
\end{quote}

The term ``halting scale'' is accepted by common usage in the
literature as the scale at which the efficiency of the cascade begins
to decline.  Reviewer B is free to challenge this usage in his or her
own papers, but we do not find its usage in this context inappropriate.

\begin{quote}
p. 14, discussion of the combined effect of beta and friction.  The
authors need to be aware that this effect has been discussed in the paper
by Galperin, Sukoriansky and Huang, ``Universal $n^{-5}$ spectrum of zonal
flows on giant planets,'' Phys. Fluids, 13, pp. 1545-1548 (2001), to which
a proper citation needs to be made.
\end{quote}

This reference was added.

\begin{quote}
p. 17, top paragraph - discussion of the width of zonal jets as
related to the steep zonal spectrum. This connection has been
explained in the same paper by Galperin \etal (2001).
\end{quote}

The paragraph mentioned is primarily concerned with the mixing scale
relevant to a passive tracer -- comments about the jet scale are made
in passing.

\begin{quote}
p. 17, 3rd paragraph from the top, discussion of the paper by Smith
and Waleffe (1999) in which energy transfer was investigated and
quartic resonances were invoked.  This discussion, almost with the
same text and the same references, can be found in a much earlier
paper by Chekhlov \etal (1996).
\end{quote}

Another citation to Chekhlov was added.

\begin{quote}
p. 27, discussion of the energy spectra on beta-plane.  The authors
describe the difference between the -5/3 and -8/3 spectra in isotropic
turbulence and then show zonal -6 spectrum on beta-plane in Fig. 7. To the
best knowledge of this reviewer, the zonal -5 spectrum has been found for
the zonal modes only in Chekhlov \etal (1996), Smith and Wallefe (1999),
and Huang \etal (2001); I am not aware of such spectrum to be found in
the entire zonal sector.  This means that the -5 spectrum, if it exists,
is delta-function like, and, thus, will retain the -5 slope even when
calculated along the zonal axis only. The authors need to check this point
very carefully.
\end{quote}

Perhaps it is a matter of definition of two-dimensional spectra,
and Reviewer B may have misunderstood our argument. To be explicit,
we meant the following.  The total energy $E$ is related to the
two-dimensional energy spectrum $\hat{E}(k_x,k_y)$ as
\[
E = \int dk_x~dk_y ~\hat{E}(k_x,k_y) = \int dk~d\theta
~k~\hat{E}(k,\theta).
\]
If we assume that the spectrum is given by
\[
\hat{E}(k,\theta) = A k^{-\alpha}\delta(\theta-\pi/2),
\]
similar to the form suggested by Reviewer B, then
\[
E = \int dk ~A k^{1-\alpha} \int d\theta ~\delta(\theta-\pi/2)
= \int dk ~A k^{1-\alpha}.
\]
Defining $\ca{E}(k)$ (the isotropic spectrum we are usually concerned
with) such that
\[
E = \int dk ~\ca{E}(k),
\]
we have that
\[
\ca{E}(k) = A k^{1-\alpha}.
\]
Hence if we predict that $\ca{E}(k) \propto k^{-5/3}$, then $\alpha =
-8/3$.  Similarly, if $\ca{E}(k) \propto k^{-5}$, then $\alpha =
-6$.

Obviously the full spectrum at least contains both parts, and in fact
is more complex altogether.  But the point is that the -8/3 and -6
spectra are just the appropriate dimensional analysis predictions for
slices of the two-dimensional spectra along the zonal and meridional
axes. We have carefully looked again at our spectra, and there are
no additional changes we wish to make in this regard. 

Thank you for your time and effort in reviewing our work.  We hope
that with these changes and comments you will now find the 
manuscript acceptable for publication in the \textit{Journal of 
Fluid Mechanics.}

\closing{Sincerely,}


\end{letter}
\end{document}


